\documentclass[11pt,a4paper,toc=bibliography,toc=listof,titlepage=firstiscover]{scrreprt}
\usepackage[utf8]{inputenc}
\usepackage{array}
\usepackage{blindtext}
\usepackage{graphicx}
\usepackage{subcaption}
\usepackage[T1]{fontenc}
\usepackage{setspace}
\usepackage{pdfsync}
\usepackage{blindtext}
\usepackage[ngerman]{babel}
\usepackage{booktabs}
\usepackage{xcolor}
\usepackage{multirow}
\usepackage{makecell}
%Schriftarten hier definieren
\usepackage{lmodern}
%\usepackage{tgheros}
%\renewcommand*\familydefault{\sfdefault}
%\usepackage{tgpagella}
\usepackage[babel,german=guillemets]{csquotes}
\usepackage[style=apa]{biblatex}
\usepackage{hyperref}
\usepackage{csquotes}
\addbibresource{res/literatur.bib}
\hypersetup{colorlinks=true, linkcolor=black, urlcolor=blue, citecolor=blue}
%Abstand der Einträge im Literaturverzeichnis als em
\setlength\bibitemsep{0.75em}
%Kein Einzug bei einem neuen Absatz
\setlength{\parindent}{0mm}
%Inhaltsverzeichnis ohne eigene Seitenzahl
\AtBeginDocument{\addtocontents{toc}{\protect\thispagestyle{empty}}} 
\usepackage[left=3cm,right=2.5cm,top=2.5cm,bottom=2cm]{geometry}

%Daten des Deckblatts
  \titlehead
  {
    {
      \begin{center}
        \vspace{2.5cm}
        \includegraphics[scale=0.3]{res/nak-logo.png}
      \end{center}
    }
  	%M.Sc. Angewandte Informatik
  }
  \subject{Hausarbeit}
  \title{Einflussfaktoren auf die Schulleistungen im Rahmen der Oberschule}
    \subtitle{MADS2100 Reporting und Visualisierung 23oB}
    \author{Leon Henne}
  \date{\small{Köln, den \today}}
  \publishers{Betreut durch Dr. Robert Stahlbock}
\begin{document}
%Titelseite erzeugen
\maketitle
%Der Nachfolgende Text wird 1,5-zeilig gesetzt
\onehalfspacing
%Römische Ziffern für Abbildungsverzeichnis
\pagenumbering{Roman}
\setcounter{page}{0}
%Inhaltsverzeichnis erzeugen
\tableofcontents
\clearpage
%Abbildungs und Tabellenverzeichnis
\listoftables
\listoffigures
\clearpage
\pagenumbering{arabic}
\setcounter{page}{1}
%Datei aus dem Ordner /chapter/ einbinden. So kann man mehrere Textdateien zu einem großen Dokument zusammenfügen.
\setcounter{chapter}{0}

\chapter[Problemstellung]{Problemstellung}

Der für jede Nation erstrebenswerte langfristige ökonomische Fortschritt wird unter anderem stark durch das vorherrschende Bildungsniveau beeinflusst (\cite[S. 1]{Cortez2008UsingDM}). 
Um dieses sich ergebende Bildungsniveau stärker zu durchleuchten und final zu verbessern, wird zur Unterstützung der Schüler und Lehrkräfte die Modellierung von Schulleistungen eingesetzt. \cite[S. 1]{Cortez2008UsingDM} 
So kann eine zeitabhängige Vorhersage der Leistungen, lernschwächere Schüler detektieren und damit Lehrkräfte frühzeitig befähigen, mit entsprechenden Maßnahmen in den Lernprozess einzugreifen. \cite[S. 2]{Namoun.2021} 
Verstärkt wurde dieser Bedarf durch die in der Vergangenheit eingetretenen Covid-Pandemie, und den damit verbundenen Schulschließungen, welche für neue erhebliche Herausforderungen sorgten. \cite[S. 2]{Clark.2021} 
Durch \cite[S. 13]{Clark.2021} konnte hierzu aufgezeigt werden, welchen positiven Effekt digitale Lernunterstützungen auf die Schülergruppen ausmachten.
Aus der Arbeit von \cite[S. 9]{Namoun.2021} geht jedoch hervor, dass bereits seit 2017 erneut das Interesse anstieg hinsichtlich der Modellierung von Lernergebnissen.
Seitdem besteht besonders ein Fokus auf die Untersuchung des Bildungsniveaus von Bachelorstudiengängen, sodass die Untersuchung weiterführender Schulen lediglich ein Anteil von in etwa 12\% besitzt. \cite[S. 11]{Namoun.2021} 
Die in den letzten Jahren erforschte Modellierung von Studierendenergebnissen lässt häufig unbeachtet, wie einzelne Faktoren innerhalb der maschinellen Lernmethoden zu den Vorhersagen führen. \cite[S. 19]{Namoun.2021}
Die Gesamtheit dieser aktuellen Gegebenheiten motiviert die nachfolgende Untersuchung des gewählten Datensatzes anhand der daran abgeleiteten Forschungsfragen.
\chapter[Zielsetzung]{Zielsetzung}

Inhalt dieses Kapitels ist die zum einen die Erläuterung des gewählten Datensatzes.
Zum Anderen werden daraus abgeleitete Forschungsfragen formuliert, welche im Rahmen dieser Arbeit durch die Entwicklung und Analyse von Visualisierungen untersucht werden. 

\section[Untersuchter Datensatz]{Datensatz}

\section[Forschungsfragen]{Forschungsfragen}
\chapter[Grundlagen]{Grundlagen}

\section[how to Visualisierungen ?]{how to Visualisierungen ?}
\section[how to analyse Visualisierungen ?]{how to analyse Visualisierungen ?}

\chapter[Untersuchung der Forschungsfrage]{Untersuchung der Forschungsfrage}

Dieses Kapitel dient der Untersuchung der Forschungsfragen. 
In Abhängigkeit des untersuchten Handlungsbereichs werden die hierfür notwendigen Datentransformationen erläutert. 
Anschließend gilt es aufgrund der Vielzahl an Attributen, die für die Untersuchung relevanten Merkmale zu identifizieren. 
Final wird die Beantwortung der Forschungsfrage, durch die Vorstellung und Analyse der entwickelten Visualisierungen vorgenommen.

\section{Individuelle und strukturelle Faktoren}

% Erläuterung der Datenstruktur: merged, left join, individual (mat\_df, por\_df)
Zur Untersuchung dieser Faktoren eignet sich besonders ein fachübergreifende Perspektive, da eine fachspezifische unterschiedliche Auswirkung nur schwierig kausal zu erklären ist.
Im Weiteren wird dadurch eine Erhöhung der Datenmenge erzielt, da alle Attribute vollständig in dem Datensatz der Portugiesisch-Schüler und Mathematik-Schüler enthalten sind.
Die Migrierung beider Datenstände wird anhand der Attribute vorgenommen, welche im, unter (\cite[]{student_performance}) verfügbaren, R-Skript benannt werden.
Insgesamt erhält der Datensatz damit 657 Einträge. 

% \item Auswahl der Variablen in Abhängigkeit zur Frage:
% \subitem individuell: Geschlecht, Alter, Besuch der Vorschule, Absicht zur Weiterbildung, außerschulische Freizeit, Alkoholkonsum unter der Woche, Alkoholkonsum am Wochenende, Gesundheitszustand
% \subitem strukturell: Grund der Schulentscheidung, Pendelzeit, häuslicher Internetzugang,
Mit der in dieser Sektion betrachteten Faktoren ergibt sich die gleichnamige Aufteilung der Attribute zu den Aspekten der individuellen Eigenschaften und Verhaltensweisen und der strukturellen Gegebenheiten.
Individuelle Merkmale lassen sich zum einen durch die Möglichkeit zur aktuellen Einflussnahme durch den Schüler selbst filtern.
Zum Anderen tritt die Absicht zur Weiterbildung zwar in diesem Kontext als Prädiktor des Lernerfolgs auf, die alleinige Absicht selbst nimmt jedoch hierauf keinen unmittelbaren Einfluss.
Damit werden die Attribute \textit{außerschulische Freizeit, Alkoholkonsum unter der Woche, Alkoholkonsum am Wochenende, Gesundheitszustand} in die Untersuchung einfließen.
Strukturelle Gegebenheiten werden im Kontext dieser Arbeit durch den Grund der Schulentscheidung, die Pendelzeit und den häuslichen Internetzugang beschrieben.
Aus diesen Merkmalen kann besonders der Pendelzeit und dem häuslichen Internetzugang ein potenzieller kausaler Zusammenhang anhand des täglichen zusätzlichen Reiseaufwands und dem Zugang zu Onlinewissen unterstellt werden.
Daher werden dazu diese beiden Attribute für die Untersuchung ausgewählt.
% \item Grafik idee:
% \subitem : (außerschulische Freizeit, dalc, walc, health) 
% \subitem Scatterplot: (Pendelzeit, Internet, G3) strukturelle Anforderungen und Noteneffekt darstellen

% ~1100 WORDS UNTIL HERE !!! 
\section{Soziales Umfeld}

\begin{itemize}
    \item Erläuterung der Datenstruktur: merged, left join, individual (mat\_df, por\_df)
    \item Auswahl der Variablen in Abhängigkeit zur Frage:
    \subitem address, famsize, Pstatus, medu, fedu, mjob, fjob, guardian, romantic, famrel, goout
    \item Grafik idee:
    \subitem Scatterplot: Bestimmen von Personengruppen anhand von visuellem (1-2) oder maschinellem (3-x) Clustering.
    \subitem Säulendiagramm: Median der Notenleistung (G3) zu den einzelnen Personengruppen
\end{itemize}

\section{Individuelle Leistungsbereitschaft}

\begin{itemize}
    \item Erläuterung der Datenstruktur: merged, left join, individual (mat\_df, por\_df)
    \item Auswahl der Variablen in Abhängigkeit zur Frage:
    \subitem studytime, failures, activities, schoolsup, famsup, paid, absences
    \item Grafik idee:
\end{itemize}
\chapter[Fazit]{Fazit}


\begin{itemize}
    \item Lernleistung lediglich durch Noten bemessen
    \item alter Datensatz
    \item Komplexität ist ok, da durch das digitale format Visualisierungen in den angezeigten Daten gefiltert werden können. (zoom, kategorie auswahl)
\end{itemize}

%Literaturverzeichnis einfügen
\renewcommand*{\UrlFont}{\rmfamily}
\printbibliography
\chapter[Anhang]{Anhang}

\begin{table}[!ht]
    \centering
    \caption{Detaillierte Darstellung des untersuchten Datensatzes}
    \label{tab:detailed_data}
    \begin{tabular}{lll}
    \hline
        \textbf{Beschreibung} & \textbf{Datentyp} & \textbf{Ausprägungen} \\ \hline \hline
        Schulbezeichnung & Binär & \makecell[l]{GP - Gabriel Pereira; \\ MS - Mousinho da Silveira} \\ \hline
        Geschlecht & Binär & F - weiblich; M - männlich \\  \hline
        Alter & Numerisch & 15 bis 22 \\  \hline
        Wohngegend & Binär & U - urban; R - ländlich \\  \hline
        Familiengröße & Binär & \makecell[l]{LE3 - kleiner oder gleich 3; \\ GT3 - mehr als 3} \\  \hline
        Zusammenleben der Eltern & Binär & \makecell[l]{T - leben gemeinsam; \\ A - leben auseinander} \\  \hline
        Mutters Bildungsgrad & Numerisch & \makecell[l]{0 - kein Bildungsgrad; \\ 1 - Grundschulabschluss (4. Klasse); \\2 – 5. bis 9. Klasse; \\3 – Oberstufenabschluss; \\4 – höherer Bildungsgrad }\\  \hline
        Vaters Bildungsgrad & Numerisch & \makecell[l]{0 - kein Bildungsgrad; \\ 1 - Grundschulabschluss (4. Klasse); \\2 – 5. bis 9. Klasse; \\3 – Oberstufenabschluss; \\4 – höherer Bildungsgrad }\\  \hline
        Mutters Arbeitsbereich & Nominal & \makecell[l]{Lehrerin; Gesundheitswesen; \\Sozialwesen (Verwaltung oder Polizei); \\zuhause; anderer Bereich }\\  \hline
        Vaters Arbeitsbereich & Nominal & \makecell[l]{Lehrer; Gesundheitswesen; \\Sozialwesen (Verwaltung oder Polizei); \\zuhause; anderer Bereich }\\  \hline
        Grund der Schulentscheidung & Nominal & \makecell[l]{Nahe dem Zuhause; Ruf der Schule; \\Kurspräferenz; anderer Grund }\\  \hline
        Erziehungsberechtigter & Nominal & Mutter; Vater; Anderer \\  \hline
        Pendelzeit & Numerisch & \makecell[l]{1 - <15 min.; 2 - 15 bis 30 min.; \\3 - 30 min. bis 1 Stunde; 4 - >1 Stunde} \\  \hline
        Lernzeit & Numerisch & \makecell[l]{1 - <2 Stunden; 2 - 2 bis 5 Stunden; \\3 - 5 bis 10 Stunden; 4 - >10 Stunden }\\  \hline
        Anzahl durchgefallener Kurse & Numerisch & 1 bis 3; 4 \\  \hline
        externe Lernunterstützung & Binär & Ja; Nein \\  \hline
        familiäre Lernunterstützung & Binär & Ja; Nein \\  \hline
        bezahlter Extraunterricht & Binär & Ja; Nein \\  \hline
        AG-Teilnahme & Binär & Ja; Nein \\  \hline
        Besuch der Vorschule & Binär & Ja; Nein \\  \hline
        Absicht zur Weiterbildung & Binär & Ja; Nein \\ \hline
        häuslicher Internetzugang & Binär & Ja; Nein \\ \hline
        Partnerliche Beziehung & Binär & Ja; Nein \\ \hline
        Qualität der Familienbeziehungen & Numerisch & von 1 - sehr schlecht bis 5 - exzellent\\ \hline
        außerschulische Freizeit & Numerisch & von 1 - sehr schlecht bis 5 - sehr gut \\ \hline
        soziale Aktivitäten & Numerisch & von 1 - sehr schlecht bis 5 - sehr gut \\ \hline
        Alkoholkonsum unter der Woche & Numerisch & von 1 - sehr schlecht bis 5 - sehr gut \\ \hline
        Alkoholkonsum am Wochenende & Numerisch & von 1 - sehr schlecht bis 5 - sehr gut \\ \hline
        Gesundheitszustand & Numerisch & von 1 - sehr schlecht bis 5 - sehr gut\\ \hline
        Fehltage & Numerisch & von 0 bis 93 \\ \hline
        erste Vorabnote & Numerisch & von 0 bis 20 \\ \hline
        zweite Vorabnote & Numerisch & von 0 bis 20 \\ \hline
        finale Note & Numerisch & von 0 bis 20 \\ \hline
    \hline
    \end{tabular}
\end{table}


\section*{Eidesstattliche Erklärung}

Hiermit erkläre ich, dass ich die vorliegende Arbeit eigenständig und ohne fremde Hilfe angefertigt habe. Textpassagen, die wörtlich oder dem Sinn nach auf Publikationen oder Vorträgen anderer Autoren beruhen, sind als solche kenntlich gemacht.

Die Arbeit wurde bisher keiner anderen Prüfungsbehörde vorgelegt und auch noch nicht veröffentlicht. \newline

Köln, den \today

\vspace{15mm}

Leon Henne


\end{document}
